\documentclass[10pt,letterpaper]{article}
\usepackage[utf8]{inputenc}
\usepackage[margin=0.75in]{geometry}
\usepackage{amsmath,amssymb}
\usepackage{booktabs}
\usepackage{enumitem}
\usepackage{hyperref}

\title{\textbf{Course Syllabus: Calculus for Biology and Social Sciences}\\
\large Based on \textit{Calculus for Biology and Medicine} (4th Ed.) by C. Neuhauser \& M. Roper}
\author{\textbf{Department of Mathematics \& Life Sciences}}
\date{}

\begin{document}

\maketitle

\section*{1. Course Overview \& Prerequisites}
\textbf{Course Description:} An introduction to differential, integral, multivariable calculus, and linear algebra designed specifically for students in biological and social sciences. Emphasis is placed on empirical data modeling, discrete/continuous dynamic systems, and stochastic processes.

\textbf{Prerequisites:} Precalculus Diagnostic Mastery including:
\begin{itemize}[noitemsep,topsep=0pt]
    \item Real numbers, intervals, inequalities, and absolute values.
    \item Exponential and logarithmic identity manipulation.
    \item Analytic geometry, straight lines, circles, and trigonometric identities.
    \item Plotting and interpreting semi-log and log-log transformed axes.
\end{itemize}

\section*{2. High-Difficulty Topics \& Core Pedagogy}
Students should dedicate additional study focus to the following mathematically challenging areas:
\begin{enumerate}[noitemsep,topsep=0pt]
    \item \textbf{Recurrence \& Autonomous Stability Analysis:} Determining stability via derivatives ($|f'(x^*)|$) and phase-line vector fields.
    \item \textbf{Multivariable Calculus \& Phase Planes:} Computing Jacobian matrices, interpreting trace-determinant criteria, and finding constrained extrema.
    \item \textbf{Integration Techniques for ODEs:} Partial fraction decomposition and two-compartment integrating factors.
\end{enumerate}

\section*{3. Module Breakdown \& Course Structure}

\begin{table}[h!]
\centering
\small
\begin{tabular}{p{2.2cm} p{3.2cm} p{5.5cm} p{5.0cm}}
\toprule
\textbf{Module \& Timeline} & \textbf{Textbook Chapters} & \textbf{Core Concepts \& Topics} & \textbf{Learning Objectives} \\
\midrule
\textbf{Module 1:} \newline Discrete Time Dynamics \newline (Weeks 1--2) 
& Ch. 1: Preview/Review \newline Ch. 2: Discrete Models \& Sequences 
& Exponential growth/decay, recurrence equations, limits of sequences, Discrete Logistic Equation, Beverton-Holt Model. 
& Formulate recursive models from verbal descriptions; evaluate limits of sequences; apply summation ($\sum$) notation. \\
\midrule
\textbf{Module 2:} \newline Differential Calculus \newline (Weeks 3--5) 
& Ch. 3: Limits \& Continuity \newline Ch. 4: Differentiation \newline Ch. 5: Applications of Derivative 
& Bisection method, Chain rule, implicit differentiation, error propagation, Mean-Value Theorem, L'H\^opital's Rule, Optimization. 
& Compute instantaneous rates of change; optimize biological single-variable equations; evaluate error bounds and limits. \\
\midrule
\textbf{Module 3:} \newline Integration \& Methods \newline (Weeks 6--8) 
& Ch. 6: Integration \newline Ch. 7: Integration Techniques 
& Riemann sums, Fundamental Theorem of Calculus, Substitution, Integration by Parts, Partial Fractions, Numerical Integration. 
& Calculate accumulated change and functional averages; solve definite/indefinite integrals required for DEs. \\
\midrule
\textbf{Module 4:} \newline Differential Equations \newline (Weeks 9--10) 
& Ch. 8: Differential Equations 
& Separable DEs, pure-time/autonomous DEs, equilibrium points \& stability, vector fields, integrating factors, compartment models. 
& Construct differential models for biological systems; determine equilibrium stability graphically and analytically. \\
\midrule
\textbf{Module 5:} \newline Linear Algebra \& Multivariable \newline (Weeks 11--13) 
& Ch. 9: Linear Algebra \newline Ch. 10: Multivariable Calculus 
& Matrices, Eigenvalues/Eigenvectors, Leslie matrices, Partial derivatives, Tangent planes, Gradient vectors, Constrained Optimization. 
& Model demographic age distributions; analyze functions of multiple variables; fit models to data using Least Squares. \\
\midrule
\textbf{Module 6:} \newline Systems of ODEs \& Probability \newline (Weeks 14--15) 
& Ch. 11: Systems of DEs \newline Ch. 12: Probability \& Statistics 
& Linear systems, nullclines, Lotka-Volterra predator-prey models, discrete/continuous probability, CLT, linear regression. 
& Classify stability of $2\times2$ systems of ODEs using eigenvalues; model stochastic biological systems and compute distributions. \\
\bottomrule
\end{tabular}
\end{table}

\newpage

\section*{4. Detailed Module Objectives \& Learning Outcomes}

\subsection*{Module 1: Discrete-Time Models, Sequences, and Difference Equations}
\begin{itemize}[noitemsep]
    \item Understand discrete-time population models including density-dependent growth models.
    \item Evaluate limits of discrete sequences and model drug absorption profiles.
\end{itemize}

\subsection*{Module 2: Single-Variable Differential Calculus \& Applications}
\begin{itemize}[noitemsep]
    \item Apply limits, continuity, and formal definitions of derivatives to biological rate systems.
    \item Perform implicit differentiation and linear approximations for error propagation.
    \item Apply optimization techniques to metabolic, physiological, and genetic systems.

\end{itemize}

\subsection*{Module 3: Integral Calculus \& Computational Techniques}
\begin{itemize}[noitemsep]
    \item Model cumulative physiological changes using the Fundamental Theorem of Calculus.
    \item Utilize substitution, integration by parts, and partial fraction decomposition to evaluate complex integrals.
\end{itemize}

\subsection*{Module 4: Differential Equations in Biology}
\begin{itemize}[noitemsep]
    \item Solve autonomous separable ordinary differential equations.
    \item Characterize stability of equilibria using phase-line analysis and vector field plots.
    \item Formulate multi-compartment pharmacological and epidemiological models.
\end{itemize}

\subsection*{Module 5: Linear Algebra, Demographic Modeling, and Multivariable Calculus}
\begin{itemize}[noitemsep]
    \item Perform matrix operations, calculate eigenvalues and eigenvectors, and project population growth using Leslie matrices.
    \item Calculate partial derivatives, linearizations, gradient vectors, and optimization of multivariate biological functions.
\end{itemize}

\subsection*{Module 6: Non-Linear Systems of Differential Equations \& Stochastic Models}
\begin{itemize}[noitemsep]
    \item Analyze phase planes, nullclines, and stability criteria for predator-prey and competitive species interaction systems.
    \item Evaluate discrete and continuous probability distributions (Binomial, Poisson, Normal, Exponential) and compute parameters via linear regression.
\end{itemize}

\end{document}